\chapter{平台定价、商业模式与大数据价格歧视}

\begin{examguide}
\textbf{本章考情}：本章是价格理论与平台商业实践的结合部，0258 高频考点（\important{专硕·核心}），"大数据杀熟"是各校论述题的常客；价格歧视的推导在学硕微观中属于必考工具（\important{学硕·热点}）。

\textbf{核心考点}：（1）三级价格歧视的 $MR_1 = MR_2 = MC$ 推导与弹性定价公式；（2）两部收费制与数量折扣的自选择设计；（3）一级价格歧视的效率与分配后果；（4）捆绑销售的利润条件与大规模捆绑定理；（5）跨期价格歧视（撇脂定价）与承诺问题；（6）行为价格歧视两期模型与个性化定价的福利分解；（7）大数据杀熟的构成要件与规制依据；（8）免费商业模式与 freemium 版本切割；（9）注意力经济与 GSP 广告竞价；（10）精准营销的福利拆解（匹配效应与定价效应）；（11）佣金归宿、最优佣金率与接入费—交易费组合。

\textbf{本章主线}：价格歧视工具箱（5.1）→ 大数据定价与规制（5.2）→ 免费与注意力（5.3）→ 佣金与抽成（5.4）。
\end{examguide}

第 4 章确立了平台定价的基本原理：平台决策的核心落在价格结构上，向谁收费、向谁补贴取决于两侧的弹性与交叉外部性强度。原理进入商业层面，就成为本章要回答的问题：平台用哪些具体工具把价格结构兑现为收入。答案集中在三组工具上，向消费者一侧实施价格歧视、向广告主一侧出售注意力、向商家一侧收取佣金抽成，它们分别对应本章的 5.1--5.2 节、5.3 节与 5.4 节。三组工具共享同一个理论底座，即价格歧视理论，本章因此从这个经典工具箱讲起。

\section{价格歧视理论}

价格歧视（price discrimination）指同一厂商对同一产品向不同消费者收取不同价格。价格歧视理论回答三个问题：什么条件下歧视可行、歧视价格如何设计、歧视的福利后果是什么。福利后果的一般结论是：\important{价格歧视改变剩余在厂商与消费者之间、以及不同消费者群体之间的分配，其效率影响取决于歧视能否扩大总产出}。对数字经济而言，这套理论从教科书走进了现实：数字服务绑定账户使跨群体转售几乎不可能，行为数据使需求差异的识别成本大幅下降，价格歧视由此成为平台定价的普遍实践，本章后续各节正是在这一背景下展开。

本节按庇古的三级分类依次展开，并在此基础上补足两类与数字产品关系最密切的歧视形态。捆绑销售把多个商品打包定价，在边际成本近零的数字产品上具有特殊威力；跨期价格歧视沿时间维度分割市场，是新品上市定价的通行做法。五个小节合起来构成平台定价的工具箱，后续三节讨论的大数据定价、免费模式与佣金抽成，都可以还原为其中某一件或某几件工具的组合运用。

\subsection{三级价格歧视}

按庇古（Pigou）的信息条件分类，价格歧视分为三级：一级按个体保留价格定价，二级通过自选择菜单定价，三级按可识别群体定价。三级价格歧视是数字平台最容易实施的形态，也是本节讨论的起点。

\begin{definition}[三级价格歧视（third-degree price discrimination）]
三级价格歧视指厂商对\textbf{可识别的不同消费者群体}收取不同价格。
\end{definition}

其前提有三：厂商具有市场势力、能够区分不同群体、转售（套利）被阻止。数字平台使第三个条件空前易于满足：数字服务的交付天然绑定账户，跨群体转售几乎不可能。在群体划分既定的条件下，厂商的最优定价规则由以下推导给出。

\begin{derivation}{三级价格歧视的最优定价}{}
设厂商把市场分割为两个可识别的群体，群体 $i$ 的需求函数为 $q_i(p_i)$（$i = 1,2$），两群体之间无法转售，生产的边际成本恒为 $c$。由于两个市场彼此隔离，厂商可以对二者独立定价，总利润等于两市场收入之和减去总生产成本：
\[
\pi = p_1 q_1(p_1) + p_2 q_2(p_2) - c\bigl(q_1 + q_2\bigr).
\]

第一步，对两个价格分别求一阶条件。价格 $p_i$ 只影响市场 $i$ 的销量，因此可以逐个市场求导。$p_i$ 提高一单位带来两项相反的效果：既有的 $q_i$ 单位销量每单位多收一元，收入增加 $q_i$；同时销量按 $dq_i/dp_i$ 减少，每减少一单位损失的边际利润为 $p_i - c$。利润最大化要求两项恰好相抵：
\[
\frac{d\pi}{dp_i} = q_i + (p_i - c)\frac{dq_i}{dp_i} = 0, \qquad i = 1,2.
\]

第二步，把一阶条件改写成弹性形式，以便看清价格由什么决定。记群体 $i$ 需求价格弹性的绝对值为 $\eta_i \equiv -\dfrac{p_i}{q_i}\dfrac{dq_i}{dp_i}$。将一阶条件两边同除以 $q_i$，再用弹性定义替换导数项，整理后得到每个市场各自的加成定价公式：
\[
p_i = \frac{\eta_i}{\eta_i - 1}\, c \quad \text{或等价地} \quad MR_1 = MR_2 = MC.
\]
两种写法含义相同。前一种说明价格是边际成本的一个加成，加成幅度 $\eta_i/(\eta_i-1)$ 完全由该市场的弹性决定，弹性越小加成越高；后一种说明厂商在两个市场之间分配产量，直到最后一单位在两处带来的边际收益相等，并且都等于边际成本。

第三步，比较两个市场的价格水平。把加成公式在 $i=1$ 与 $i=2$ 时相除，边际成本 $c$ 被约去，价格比只剩下弹性：
\[
\frac{p_1}{p_2} = \frac{\eta_1(\eta_2 - 1)}{\eta_2(\eta_1 - 1)}.
\]
注意到加成函数 $\eta/(\eta-1) = 1 + 1/(\eta-1)$ 关于 $\eta$ 严格递减，故给定 $\eta_2$ 时，群体 1 的弹性 $\eta_1$ 越小，比值 $p_1/p_2$ 越大。这就得到三级价格歧视的基本结论：\important{需求弹性更小的群体被收取更高价格}。
\end{derivation}

推导的经济含义可以概括为一句话：价格差异来自弹性差异，经由利润最大化条件传导为定价规则。弹性小的群体对价格不敏感，涨价带来的销量损失有限，厂商向其转嫁更多；弹性大的群体一涨价就流失，厂商只能让利留客。学生票、老年票、软件的"学生版/专业版"定价都循此理：学生支付能力低、替代选择多，需求弹性大，因而拿到低价。\mn{三级价格歧视计算题三步法：第一步由 $MR_1 = MR_2 = MC$ 解出两市场产量；第二步把产量代回各自需求函数读出价格；第三步点明"弹性小的市场价高"的结论句。}三级价格歧视的可行性条件也由此清楚：厂商须具有市场势力（否则价格由市场决定）、能够低成本识别群体（否则分割无从实施）、能够阻止转售（否则低价群体会把商品倒卖给高价群体，价差被套利抹平）。

与统一定价相比，三级价格歧视的福利后果是分层的。高弹性群体获益，其价格从统一价下降；低弹性群体受损，其价格上升；厂商利润必然上升，因为统一定价始终是歧视定价集合中的一个可行选项，厂商选择歧视意味着歧视带来的利润更高。总福利的变化方向则不确定，取决于歧视是否使市场总产出扩大：若歧视只是把既有销量在两群体之间重新分配，总产出不变而配置扭曲加深，总福利下降；若歧视使原本被统一高价挡在市场之外的高弹性群体开始购买，总产出扩大，总福利可能上升。\important{产出是否扩大是价格歧视福利评价的一般判据}，5.2 节的个性化定价与 5.3 节的精准营销都要回到这一条。

三级价格歧视在数字服务中的实施成本极低，这是它成为平台默认工具的直接原因。视频平台的会员分层（学生认证优惠、连续包月与单月购买的差价）、云服务商面向高校与初创企业的专项折扣、在线教育按地区差异化的课程定价，都以可核验的身份标签划分群体，识别成本接近于零。中国网络视频用户规模超过 10 亿（中国互联网络信息中心，2023），会员分层因此覆盖极大的人群基数，微小的价差在总量上就是可观的收入差异。转售条件的变化同样关键：传统行业受制于套利，价格歧视只能依托难以转售的服务型商品（电影票、景区门票、理发服务）；数字服务的交付绑定账户，播放权限、下载额度、API 调用配额都随账户走，转售在技术上被直接封死，三级价格歧视的适用范围由此从少数行业扩展到几乎全部线上业务。

\begin{figure}[H]
\centering
\begin{tikzpicture}[>=Stealth, scale=0.9]
% 左图：市场1（弹性大）
\begin{scope}[shift={(0,0)}]
% 坐标轴
\draw[->, gray] (0,0) -- (5.2,0);
\draw[->, gray] (0,0) -- (0,5.2);
% D1: p = 4 - q；MR1: p = 4 - 2q；MC = 1
\draw[thick, second] (0,4) -- (4,0);
\draw[thick, main] (0,4) -- (2,0);
\draw[dashed, thick, winered] (0,1) -- (4.4,1);
\node[font=\small, color=winered, anchor=south] at (4.4,1) {$MC$};
\node[font=\small, color=second, anchor=west] at (0.1,4.35) {$D_1$（弹性大）};
\node[font=\small, color=main, anchor=west] at (2.05,0.4) {$MR_1$};
% 均衡：MR1 = MC → q1* = 1.5；投影到 D1 → p1* = 2.5
\fill[second] (1.5,1) circle (2.5pt);
\fill[second] (1.5,2.5) circle (2.5pt);
\draw[dashed, gray] (1.5,1) -- (1.5,0) node[anchor=north, font=\small] {$q_1^*$};
\draw[dashed, gray] (1.5,1) -- (1.5,2.5);
\draw[dashed, gray] (1.5,2.5) -- (0,2.5) node[anchor=east, font=\small] {$p_1^*$};
\end{scope}
% 右图：市场2（弹性小）
\begin{scope}[shift={(6.6,0)}]
\draw[->, gray] (0,0) -- (5.6,0);
\draw[->, gray] (0,0) -- (0,5.2);
% D2: p = 5 - 0.8q（陡峭，弹性小）；MR2: p = 5 - 1.6q；MC = 1
\draw[thick, second] (0,5) -- (5,1);
\draw[thick, main] (0,5) -- (2.5,1);
\draw[dashed, thick, winered] (0,1) -- (4.4,1);
\node[font=\small, color=winered, anchor=south] at (4.0,1) {$MC$};
\node[font=\small, color=second, anchor=south] at (3.95,1.6) {$D_2$（弹性小）};
\node[font=\small, color=main, anchor=west] at (1.05,1.52) {$MR_2$};
% 均衡：MR2 = MC → q2* = 2.5；投影到 D2 → p2* = 5 - 0.8×2.5 = 3
\fill[second] (2.5,1) circle (2.5pt);
\fill[second] (2.5,3) circle (2.5pt);
\draw[dashed, gray] (2.5,1) -- (2.5,0) node[anchor=north, font=\small] {$q_2^*$};
\draw[dashed, gray] (2.5,1) -- (2.5,3);
\draw[dashed, gray] (2.5,3) -- (0,3) node[anchor=east, font=\small] {$p_2^*$};
\end{scope}
\end{tikzpicture}
\caption{三级价格歧视}
\label{fig:pricedisc}
\end{figure}

图~\ref{fig:pricedisc} 直观呈现了三级价格歧视的定价结构。两个市场共享同一条 $MC$ 线，需求曲线的斜率却不同：市场 1 的需求曲线平缓，对应弹性大；市场 2 陡峭，对应弹性小。读图的顺序与推导的顺序完全一致：均衡条件 $MR_1 = MR_2 = MC$ 先在每个市场确定最优产量，即各自 $MR$ 曲线与 $MC$ 线的交点；再由需求曲线在该产量处向上投影，读出对应的价格。结果是市场 1 的最优价格 $p_1^* = 2.5$ 低于市场 2 的 $p_2^* = 3$，弹性小的市场价格更高这一结论在图上一目了然。

\subsection{二级价格歧视与两部收费}

当厂商无法识别消费者类型时，三级价格歧视无从实施。替代思路是让消费者自我暴露类型：设计一组"菜单"，不同类型的消费者会自发选择不同选项，数量折扣、版本差异化皆是如此。

\begin{definition}[二级价格歧视（second-degree price discrimination）]
二级价格歧视指厂商无法识别消费者个体类型时，通过设计自选择菜单（数量折扣、版本差异）让消费者自我暴露类型的定价方式。
\end{definition}

自选择菜单的经典实现是两部收费制（two-part tariff）：一笔与消费量无关的固定入场费，加上按消费量计算的单价，消费者的总支出为 $T = A + pq$。会员制、月租加通话费都是这一结构的现实形态。以下先在消费者同质的简化情形下推导其最优设计，再讨论类型异质带来的复杂性。

\begin{derivation}{两部收费制的最优设计}{}
设消费者需求为 $q(p) = a - p$，厂商边际成本为 $c$，且所有消费者的需求函数相同（先考察同质情形）。两部收费制下厂商的利润由两块构成，一块是入场费收入 $A$，另一块是按量销售的边际利润 $(p-c)q(p)$：
\[
\pi = A + (p - c) q(p).
\]

第一步，确定给定单价 $p$ 时厂商能够收取的最高入场费。消费者是否参与取决于净剩余是否非负：支付入场费之后，消费本身带来的消费者剩余必须至少覆盖这笔入场费。因此厂商可榨取的入场费上限恰好等于该单价下的全部消费者剩余。反需求曲线 $p = a - q$ 给出第 $q$ 个单位买者的支付意愿，消费者剩余等于支付意愿超出实付价格的部分沿全部购买量累积：
\[
A^* = CS(p) = \int_p^a (a - t)\, dt = \frac{(a-p)^2}{2}.
\]

第二步，把入场费代回利润函数，消去 $A$，得到只含单价 $p$ 的目标函数：
\[
\pi(p) = \frac{(a-p)^2}{2} + (p-c)(a-p).
\]
式中第一项是入场费收入，它随 $p$ 上升而下降，因为单价越高消费者剩余越少、可榨取的入场费越低；第二项是按量销售的边际利润，在 $p$ 高于 $c$ 时为正。两项随单价变动的方向相反，最优单价由二者的边际权衡决定。

第三步，对 $p$ 求导求解。第一项求导得 $-(a-p)$，第二项按乘积法则求导得 $(a-p)-(p-c)$，两者合并后前后项相互抵消，只剩一项：
\[
\frac{d\pi}{dp} = -(a-p) + (a-p) - (p-c) = -(p-c) = 0 \;\Rightarrow\; p^* = c.
\]
导数的形式本身就说明了问题：单价每提高一元，从按量销售中多赚的钱恰好被入场费的损失抵消，唯一残留的 $-(p-c)$ 代表单价高于边际成本所造成的销量扭曲。只有当 $p = c$、扭曲归零时利润才达到最大。将 $p^*=c$ 代回入场费公式，得 $A^* = (a-c)^2/2$，即边际成本定价下的消费者剩余。最优两部收费由此确定：\important{单价定在边际成本，入场费榨取全部消费者剩余}。
\end{derivation}

这一结果把两部收费的双重功能分得很清楚：单价负责效率，入场费负责分配。$p^* = c$ 意味着定价规则与完全竞争一致，每一个支付意愿高于边际成本的单位都被生产出来，无谓损失为零；$A^* = CS(c)$ 意味着这些效率收益全部落入厂商口袋，消费者的净剩余恰好为零。\mn{两部收费速记：单价 $p^*=c$ 管效率、入场费 $A^*=CS(c)$ 管分配；同质消费者下厂商利润等于一级价格歧视，但只需知道需求曲线形状，无需识别任何个体。}厂商在此实现了与一级价格歧视相同的利润，所需信息却少得多：只要知道需求曲线的形状即可，无需识别每一个消费者的具体身份。信息门槛低而榨取能力强，这正是两部收费在实务中被广泛采用的原因。

会员费加低价商品的组合，是两部收费最直白的现实版本。山姆会员店、开市客收取年费，商品以接近成本的价格销售，利润主要来自会员费，商品加价的贡献很小；电信运营商的月租加通话费、健身房的年卡加私教课时费同样属于这一结构。数字平台把它用得更彻底：电商平台的付费会员以年费换取免运费与专属折扣，外卖平台的会员卡以月费换取免配送费，逻辑都是先用入场费一次性锁定剩余，再把边际消费的价格压到尽可能低以刺激使用频次（第 15 章从企业经营视角再次讨论这一工具的运用）。云计算的定价也有类似结构，基础套餐费对应入场费，超出部分按实际用量计价。

消费者类型异质时，单一的两部收费不再最优。若入场费按高需求类型的剩余设定，低需求类型将因净剩余为负而退出市场；若按低需求类型设定，高需求类型将获得正剩余，榨取不完全。厂商的解法是提供一组套餐让消费者自选，这就回到了二级价格歧视的本义。最优菜单须满足自选择约束，即每一类型选择为其设计的套餐所得剩余不低于选择其他套餐，其典型结果是：高需求类型得到接近有效率的消费量但被榨取大部分剩余，低需求类型的消费量被人为压低，以此削弱高类型冒充低类型的诱惑（第 2 章信息甄别模型中的保险合同设计是同一逻辑的另一应用）。

数量折扣是自选择菜单最常见的形态。厂商设定随购买量递减的平均价格，例如前 10 个单位每单位 5 元、此后每单位 3 元，高需求消费者自然买得多、享受较低的平均价，低需求消费者买得少、承担较高的平均价，群体划分由购买行为自动完成，厂商无需事先识别任何人。数字业务中的阶梯计价随处可见：云服务的存储与流量按使用量分档，调用量越大单价越低；应用程序接口（API）服务按月调用次数设置价格档位；软件即服务（SaaS）按坐席数分层报价，团队规模越大人均价格越低。这些价目表在形式上是折扣，在经济学上则是把两部收费的思想推广到连续的消费量维度：每一档的档位费相当于一次小额入场费，档内单价则向边际成本靠拢。

\subsection{一级价格歧视}

三级、二级价格歧视都只利用了部分支付意愿信息，最极端的情形则是厂商掌握每个消费者对每一单位的支付意愿，此时的价格歧视被称为一级价格歧视。

\begin{definition}[一级价格歧视（first-degree price discrimination）]
一级价格歧视，又称完全价格歧视（perfect price discrimination），指厂商对每一单位商品按该单位买者的保留价格收费，消费者剩余被完全榨取的定价方式。
\end{definition}

此时需求曲线上的每一点都被单独定价。这一形态看似纯理论，其效率性质却为理解数字时代的个性化定价提供了基准。

\begin{derivation}{完全价格歧视的产出与剩余分配}{}
设市场的反需求函数为 $p(q)$，厂商的边际成本函数为 $MC(q)$。完全价格歧视下，厂商对第 $q$ 个单位收取的价格恰好等于该单位买者的保留价格 $p(q)$。

第一步，确定边际收益曲线的位置。统一定价时，厂商多卖一单位必须对全部销量降价，边际收益为 $MR = p(q) + q\,\dfrac{dp}{dq}$，其中第二项为负，代表降价对既有销量造成的收入侵蚀，故 $MR < p(q)$。完全价格歧视下这一侵蚀不再存在：新增的那一单位以其自身的保留价格单独成交，既有单位的价格丝毫不受影响。边际收益曲线因此与需求曲线重合：
\[
MR(q) = p(q).
\]

第二步，求最优产量。利润最大化条件仍为 $MR = MC$，把上式代入即得
\[
p(q^*) = MC(q^*).
\]
这恰好是完全竞争市场的均衡条件，故 $q^* = q_c$。厂商之所以愿意一直卖到价格等于边际成本，是因为多卖一单位的代价仅为该单位的生产成本，不必再为此对其他单位让价。

第三步，计算剩余的分配。厂商的总收入等于每一单位保留价格的累积，即需求曲线下方直到 $q^*$ 的全部面积 $\int_0^{q^*} p(q)\,dq$；总成本为 $\int_0^{q^*} MC(q)\,dq$。二者之差为
\[
\pi = \int_0^{q^*}\bigl[p(q) - MC(q)\bigr]dq,
\]
这一积分正是需求曲线与边际成本曲线之间的全部面积，也就是社会总剩余。由于每位消费者在每一单位上的支付意愿都被价格取尽，总消费者剩余为零。与完全竞争相比：\important{效率相同（无谓损失为零），分配全然不同（剩余全部归厂商）}。
\end{derivation}

这一结果把效率与分配两个维度彻底分开，其理论意义超出价格歧视本身。教科书通常以无谓损失作为批评垄断的理由，完全价格歧视却表明，无谓损失源于统一定价这一约束，并非源于垄断地位本身：一旦垄断者能够对每一单位单独定价，它反而会把产量扩张到有效率的水平，因为多卖一单位不再需要对既有销量让价。福利评价的焦点因此从总量转向分配：总剩余达到最大，消费者却一无所获。这也解释了为什么监管对个性化定价的关切集中在公平与分配，而较少诉诸效率损失。

一级价格歧视长期被视为理论上的极端情形，因为厂商无从知晓每个消费者的保留价格。数字技术正在削弱这一信息约束：平台掌握高度个体化的需求信息（浏览与搜索历史、支付记录、页面停留时长、设备型号与操作系统），机器学习模型据此估计个体的价格敏感度，报价在技术上已能做到千人千面。\mn{三类价格歧视的判别速记：三级=按\textbf{群体}定价（可识别群体），二级=按\textbf{数量/版本}定价（自选择），一级=按\textbf{个体保留价格}定价（完全信息）。大数据杀熟接近一级价格歧视。}现实中的个性化定价与完全价格歧视仍有距离，估计误差、消费者的比价与投诉、监管约束都限制着榨取的完全程度，但方向是明确的。教科书假设正在变成监管议题，这就是 5.2 节讨论大数据杀熟的理论位置。

\subsection{捆绑销售}

前面三种价格歧视都在单一商品内部分割市场，捆绑销售则换了一个维度：把多个商品打包成一个组合定价，利用消费者对不同商品估价的差异结构来榨取剩余。数字产品的复制边际成本近乎为零，打包几乎不增加成本，捆绑因此成为数字经济中威力最大的定价工具之一（第 1 章讨论零边际成本行业的定价机制时曾预告这一主题）。

\begin{definition}[捆绑销售（bundling）]
捆绑销售指厂商把两种或多种商品组合成一个整体定价出售的行为。只提供组合、不允许单独购买的称为\textbf{纯捆绑}（pure bundling），既提供组合又允许单独购买的称为\textbf{混合捆绑}（mixed bundling）。
\end{definition}

捆绑为什么能提高利润？关键在于打包会改变消费者估价的分布形态。以下先用一个两商品两消费者的算例说明机制，再把结论推广到多商品的一般情形。

\begin{derivation}{捆绑的利润条件与大规模捆绑}{}
\textbf{第一部分：两商品算例。}设厂商生产 A、B 两种商品，边际成本均为零，市场上有两位消费者。消费者甲对 A 的支付意愿为 10、对 B 为 4；消费者乙对 A 为 4、对 B 为 10。两人的估价总和相同，但方向相反。

分别定价时，厂商对每种商品单独选择价格。以商品 A 为例：定价 10 时只有甲购买，收入 10；定价 4 时两人都买，收入 8。故 A 的最优价格为 10、利润 10。商品 B 由对称性同样得利润 10。分别定价的总利润为
\[
\pi_{\text{分售}} = 10 + 10 = 20.
\]

纯捆绑时，厂商只出售 A 与 B 的组合包。甲对组合包的支付意愿为 $10 + 4 = 14$，乙为 $4 + 10 = 14$，两人完全一致。厂商定价 14 时两人都购买：
\[
\pi_{\text{捆绑}} = 14 \times 2 = 28 > 20.
\]

利润提高的来源清晰可见。单独出售时，两位消费者对同一商品的估价相差 6，厂商只能二选一：要么高价成交而放弃一半市场，要么低价成交而放弃高估价者的剩余。打包之后，两人在 A 上的估价差与在 B 上的估价差方向相反，相互抵消，组合包的估价分布从分散变为集中，厂商用一个价格就同时取尽了两人的全部支付意愿。由此得到捆绑的利润条件：\important{消费者对不同商品的估价负相关（或至少不完全正相关）时，捆绑提高利润}。负相关越强，打包后的估价越集中，可榨取的剩余越多；若估价完全正相关，打包不改变估价的分散程度，捆绑无利可图。

\textbf{第二部分：大规模捆绑。}设厂商拥有 $n$ 种商品，边际成本均为零，消费者对第 $i$ 种商品的估价 $v_i$ 相互独立、同分布，均值为 $\mu$、方差为 $\sigma^2$。消费者对包含全部 $n$ 种商品的组合包的支付意愿为 $S_n = \sum_{i=1}^{n} v_i$，人均估价为 $S_n/n$。由估价之间的独立性，方差可以逐项相加，人均估价的方差为
\[
\operatorname{Var}\Bigl(\frac{S_n}{n}\Bigr) = \frac{n\sigma^2}{n^2} = \frac{\sigma^2}{n} \xrightarrow{\;n \to \infty\;} 0.
\]
方差趋于零意味着：随着捆绑规模扩大，所有消费者对组合包的人均估价都向同一数值 $\mu$ 收敛，个体之间的偏好差异被平均掉。厂商据此将组合包定价为 $n(\mu - \varepsilon)$，其中 $\varepsilon$ 为任意小的正数。当 $n$ 足够大时，人均估价低于 $\mu - \varepsilon$ 的消费者比例趋于零，几乎全部消费者都会购买，厂商的人均收入趋于 $\mu$。而 $\mu$ 正是消费者对单件商品的平均支付意愿，也就是完全价格歧视所能榨取的人均上限。结论是：\important{商品种类足够多且估价相互独立时，大规模捆绑的利润逼近完全价格歧视}。
\end{derivation}

大规模捆绑的结论解释了数字产品定价的一个显著特征：种类越多，打包越有利。\mn{大数定律：相互独立同分布的随机变量，其样本均值随样本量增大向总体均值收敛，即 $\frac{1}{n}\sum_{i=1}^{n}v_i \to \mu$。直观含义是个体差异被平均掉，这是大规模捆绑使估价分布收窄的数学依据。}软件套件把文字处理、表格、演示、邮件与云存储装进一个订阅，视频平台把电影、剧集、综艺与体育赛事装进一个会员，音乐平台把数千万首曲目装进一个包月，云厂商把计算、存储、数据库与安全服务装进一个资源包。每位用户真正高频使用的可能只有其中三五项，但每个人高频使用的那三五项各不相同，估价的独立性恰好满足定理的前提，打包因此能用一个价格同时留住偏好差异极大的人群。

捆绑的成本条件在数字产品上格外有利。实体商品的捆绑受制于生产与库存成本，多装一件就多一份边际成本，组合包的规模很快触及成本约束；数字产品的复制边际成本近乎为零，把第 1000 部影片加入会员目录几乎不增加成本，却能覆盖一批专为这部影片付费的用户。这一成本上的不对称使数字平台的最优捆绑规模远大于传统零售，也解释了内容平台持续扩充片库、软件厂商不断把新功能并入订阅的行为逻辑。

捆绑同时是竞争政策关注的行为。当厂商在某一商品上具有市场势力，并通过捆绑把这一势力传导到相邻的竞争性商品上时，捆绑就可能构成搭售（tying），排挤相邻市场的竞争者。定价工具与排他行为的边界由第 6 章的竞争政策处理，其判断关键在于捆绑是否产生了排除竞争的效果，而非捆绑行为本身。

\subsection{跨期价格歧视与撇脂定价}

以上几种歧视都在同一时点分割市场，跨期价格歧视则沿时间维度分割：先以高价卖给最急迫的消费者，再降价覆盖愿意等待的消费者。新品上市的价格轨迹、耐用品的逐年降价、演唱会门票的分批发售都属于这一形态。

\begin{definition}[撇脂定价（skimming pricing）]
撇脂定价指厂商在新产品上市初期定高价，先榨取支付意愿最高的消费者的剩余，随后逐步降价以覆盖支付意愿较低的消费者的跨期定价策略。因其形似逐层撇取牛奶表面的浮脂而得名。
\end{definition}

跨期歧视的实施条件与前几种不同：它不要求厂商识别消费者身份，只要求消费者在等待意愿上存在差异。但它引入了一个前几种歧视没有的难题：消费者会预期到未来的降价，从而策略性地推迟购买。以下模型刻画厂商与消费者之间的这一互动。

\begin{derivation}{两期撇脂定价与承诺问题}{}
设消费者对产品的估价 $v$ 在 $[0,1]$ 上均匀分布，人口总量标准化为 1，厂商边际成本为零，销售分两期进行。消费者与厂商的折现因子同为 $\delta \in [0,1)$，$\delta$ 越大表示消费者越有耐心、越不介意等待。厂商在第一期定价 $p_1$、第二期定价 $p_2$，且无法在第一期承诺第二期的价格。

第一步，确定第一期的购买者范围。每位消费者面对两个选择：第一期就买，净得 $v - p_1$；等到第二期再买，净得 $\delta(v - p_2)$。存在一个临界估价 $\hat v$ 使两个选择无差异：
\[
\hat v - p_1 = \delta(\hat v - p_2).
\]
估价高于 $\hat v$ 的消费者急于消费，第一期购买；估价低于 $\hat v$ 的消费者选择等待。把上式整理，得到第一期价格与临界值之间的关系：
\[
p_1 = \hat v(1-\delta) + \delta p_2.
\]

第二步，求第二期的最优价格。第一期结束后，市场上剩下的是估价低于 $\hat v$ 的消费者，其估价在 $[0,\hat v]$ 上均匀分布。厂商面对这一剩余需求重新定价：定价 $p_2$ 时购买者为估价落在 $[p_2,\hat v]$ 区间的消费者，人数为 $\hat v - p_2$，第二期收入为 $p_2(\hat v - p_2)$。对 $p_2$ 求导并令其为零，得 $\hat v - 2p_2 = 0$，故
\[
p_2^* = \frac{\hat v}{2}, \qquad \pi_2^* = \frac{\hat v}{2}\Bigl(\hat v - \frac{\hat v}{2}\Bigr) = \frac{\hat v^2}{4}.
\]
这一步是厂商在第二期的最优反应，它不受第一期任何承诺的约束，消费者也正是预见到了这一点。

第三步，求第一期的最优定价。厂商选择 $p_1$ 等价于选择临界值 $\hat v$。把 $p_2^* = \hat v/2$ 代回第一步的关系式：
\[
p_1 = \hat v(1-\delta) + \delta\frac{\hat v}{2} = \hat v\Bigl(1 - \frac{\delta}{2}\Bigr).
\]
两期总利润等于第一期收入加上贴现后的第二期利润，第一期购买者人数为 $1 - \hat v$：
\[
\pi = \hat v\Bigl(1-\frac{\delta}{2}\Bigr)(1-\hat v) + \frac{\delta \hat v^2}{4}.
\]
展开并把两个 $\hat v^2$ 项合并：
\[
\pi(\hat v) = \Bigl(1-\frac{\delta}{2}\Bigr)\hat v - \Bigl(1 - \frac{3\delta}{4}\Bigr)\hat v^2.
\]
对 $\hat v$ 求导并令其为零：
\[
\Bigl(1-\frac{\delta}{2}\Bigr) - 2\Bigl(1-\frac{3\delta}{4}\Bigr)\hat v = 0 \;\Rightarrow\; \hat v^* = \frac{2-\delta}{4-3\delta}.
\]
代回前两式，得两期的最优价格
\[
p_1^* = \frac{(2-\delta)^2}{2(4-3\delta)}, \qquad p_2^* = \frac{2-\delta}{2(4-3\delta)}.
\]
两者之比为 $p_1^*/p_2^* = 2-\delta > 1$（因 $\delta < 1$），故\important{第一期价格高于第二期价格}，价格轨迹先高后低，这正是撇脂定价的形态。

第四步，与单期垄断作比较。若厂商只销售一期，面对 $U[0,1]$ 的需求，最优价格为 $1/2$、利润为 $1/4$。现取 $\delta = 1/2$ 代入上面的解：$\hat v^* = 0.6$、$p_1^* = 0.45$、$p_2^* = 0.3$，两期总利润为
\[
\pi = 0.45 \times 0.4 + 0.5 \times \frac{0.36}{4} = 0.18 + 0.045 = 0.225.
\]
多了一件定价工具，利润却\important{低于单期垄断的 $0.25$}。
\end{derivation}

最后一步的比较结果值得停下来想清楚。厂商掌握了更多定价工具（两个价格而非一个），利润反而下降，原因在于承诺能力的缺失。消费者预见到第二期必然降价，高估价者中的一部分选择等待，第一期需求随之削弱；厂商为了留住这些人只能压低第一期价格，第一期定价因此被消费者的等待预期"锚"住。工具的增加没有转化为利润，反而制造出厂商与自己未来行为之间的一场博弈。这一现象在耐用品垄断文献中被称为科斯猜想（Coase conjecture）：若消费者足够有耐心、厂商降价可以足够频繁，垄断价格将迅速跌向边际成本，垄断力量被自身的降价能力瓦解。

现实中的撇脂定价之所以仍然有效，靠的正是各种恢复承诺能力的安排。限量发售与停产承诺让消费者相信等待未必有货；产品迭代把降价与新品发布捆在一起，等待者最终拿到的是过时型号；租赁替代出售可以让厂商避开与未来的自己竞争。数字产品还有一类特殊手段，即版本控制：首发版本与后续促销版本在内容或服务上作出区隔（首发限定内容、抢先体验、赛季通行证），使等待的代价从"晚一点用"变成"拿不到同样的东西"。

撇脂定价在数字业务中的典型场景包括智能手机与消费电子的新品定价（发布时定高价，次年新品上市时旧款正式降价）、游戏的首发定价与后续折扣、在线课程与知识付费产品的首发优惠、流媒体内容的窗口期发行（院线、付费点播、会员免费依次开放）。\mn{撇脂与渗透的对照速记：撇脂＝先高后低，适用于估价分散、无网络效应的产品；渗透＝先低后高，适用于网络效应强、需要跨越临界规模的产品。"为什么社交产品免费而游戏主机首发高价"，答案就在这一对比里。}值得注意的是游戏行业的季节性促销已形成稳定预期，反过来训练出一批专门等待打折的玩家，这正是模型中 $\delta$ 较大的消费者，厂商的应对是把首发与限定内容绑定以抬高等待成本。与撇脂相对的策略是渗透定价（penetration pricing）：新品以低价甚至免费快速积累用户规模，待网络效应形成后再提价或另寻变现渠道。两者的选择取决于产品是否具有网络外部性：网络效应强的产品（社交、平台类）用户规模本身创造价值，渗透定价占优（第 3 章的临界规模与第 4 章的补贴定价给出了完整分析）；网络效应弱、消费者估价分散的产品（内容、硬件、工具软件）则撇脂定价占优。

\section{大数据与个性化定价}

5.1 节的价格歧视工具都以信息为前提：三级歧视要求识别群体，二级歧视要求设计出能让人自我暴露类型的菜单，一级歧视要求知道每个人的保留价格。信息约束越紧，厂商能够实施的歧视级别越低。数据技术改变的正是这一约束：平台在提供服务的过程中持续积累行为记录，识别成本随之下降，价格歧视的可行边界向一级歧视的方向推进。

本节沿着这一推进展开四个层次的讨论：行为数据如何转化为定价能力（5.2.1 节）、个性化定价如何重塑福利在平台与消费者之间的分配（5.2.2 节）、数据用于匹配时的福利效果与用于定价时有何区别（5.2.3 节），以及法律与监管如何回应（5.2.4 节）。前三个层次分别对应一个推导，最后一个层次落到中国现行的规制框架。

\subsection{行为价格歧视}

三级价格歧视的前提是厂商能够识别群体，大数据时代的新变化在于：平台无须事先"认识"消费者，只要积累消费历史，就能从行为本身推断支付意愿。"大数据杀熟"正是这一机制的社会标签，其经济学名称是行为价格歧视。

\begin{definition}[行为价格歧视（behavior-based price discrimination）]
行为价格歧视是指厂商利用消费者的\textbf{历史行为数据}推断其支付意愿并据此差异化定价的行为；"大数据杀熟"是其公众称谓：老用户看到的价格高于新用户，活跃用户高于偶发用户。
\end{definition}

为什么历史行为能成为定价依据？关键在于购买行为本身携带信息：买与不买的决策揭示了保留价格的高低。刻画这一机制需要动态框架，即两期模型：第一期的购买决策揭示类型，第二期据此差异化定价。

\begin{derivation}{两期行为定价：老客户为何被加价}{}
设消费者对平台服务的保留价格为 $v$，$v$ 在 $[0,1]$ 上均匀分布，人口总量标准化为 1，平台的边际成本为零。销售分两期进行，第一期平台尚无任何个体信息，只能统一定价 $p_1$。

第一步，识别购买行为所携带的信息。第一期定价 $p_1$ 时，购买者必然满足 $v \geq p_1$，未购买者必然满足 $v < p_1$。购买这一行为本身因此把人群切成了两半：平台虽然仍不知道任何人的确切估价，却知道了每个人的估价落在哪个区间。用条件分布表述，第一期购买者（"老客户"）的估价服从 $U[p_1, 1]$，未购买者（"新客户"）的估价服从 $U[0, p_1]$。这就是行为数据转化为定价能力的全部机制：\important{买与不买本身是一个关于支付意愿的自然信号}。

第二步，求老客户的第二期最优价格。平台面对估价服从 $U[p_1,1]$ 的老客户群体定价 $p_L$。均匀分布下，愿意支付 $p_L$ 的人数占老客户总数的比例，等于区间 $[p_L,1]$ 的长度与区间 $[p_1,1]$ 长度之比，故需求为 $q_L = (1-p_L)/(1-p_1)$，从老客户处获得的第二期利润为
\[
\pi_L = \frac{p_L(1 - p_L)}{1-p_1}.
\]
分母 $1-p_1$ 与 $p_L$ 无关，最大化 $\pi_L$ 因而等价于最大化分子 $p_L(1-p_L)$，其顶点在 $p_L = 1/2$。但此处存在一个约束：老客户的估价都不低于 $p_1$，定价一旦低于 $p_1$，全体老客户都会购买，继续降价只是白白让利。因此当 $p_1 > 1/2$ 时，目标函数在可行区间 $[p_1,1]$ 上已处于递减段，最优解落在左端点。两种情形合并写作
\[
p_L^* = \max\Bigl\{p_1,\; \frac{1}{2}\Bigr\}.
\]

第三步，求新客户的第二期最优价格。新客户的估价服从 $U[0,p_1]$，定价 $p_N$ 时购买者占新客户的比例为 $(p_1 - p_N)/p_1$，利润为 $\pi_N = p_N(p_1-p_N)/p_1$。同样地，分母与 $p_N$ 无关，对分子求导得 $p_1 - 2p_N = 0$，故
\[
p_N^* = \frac{p_1}{2}.
\]

第四步，比较两个价格。当 $p_1 \leq 1/2$ 时，$p_L^* = 1/2$ 而 $p_N^* = p_1/2 \leq 1/4$；当 $p_1 > 1/2$ 时，$p_L^* = p_1$ 而 $p_N^* = p_1/2$。两种情形下都有
\[
p_L^* > p_N^*,
\]
故\important{老客户被收取更高价格}。同时 $p_L^* \geq p_1$ 恒成立，即老客户在第二期面对的价格不低于第一期：第一期的购买行为暴露了较高的支付意愿，平台随即据此抬价。
\end{derivation}

模型的结论看起来对平台完全有利，但把消费者的预期纳入进来，图景就复杂了。理性的消费者会预见到"第一期购买将导致第二期被加价"这一后果，购买的实际代价因此高于标价 $p_1$，其中还包含了被识别为高支付意愿者的隐性成本。第一期的购买激励由此下降，需求被抑制。平台面临的权衡因而十分清楚：杀熟的收益来自第二期对老客户的剩余榨取，代价是第一期的需求收缩与用户积累放缓。第一期定价越低，愿意暴露自己的用户越多，第二期可供榨取的对象也越多，这解释了平台为何常以补贴价换取用户的首次购买。

对消费者福利的影响则是双向的。第二期价格上升直接损害老客户，而平台为吸引"未来的数据源"所压低的第一期价格，又对早期消费者构成补偿。两种效应的净方向取决于消费者的耐心程度、平台侧的竞争强度，以及消费者能否隐藏身份（更换账号、清除记录、使用隐私模式）。行为价格歧视在理论上难有统一的福利结论，原因正在于此；考试中若要求评价其福利后果，应当分情形讨论，避免给出单一判断。

从价格歧视的分类体系看，行为价格歧视介于三级与一级之间。它的群体划分内生于购买行为，由消费者自己的选择产生，无需依赖年龄、职业这类先验的外生标签；它的价格随个体消费历史逐步向该消费者的保留价格靠拢，比按群体统一制定价格的三级歧视精细得多。\mn{两期模型大白话：第一期低价"钓"出高支付意愿用户并积累其消费数据，第二期对老客户加价"收割"。结论速记：老客户价 $p_L = \max\{p_1,\; 1/2\}$、新客户价 $p_N = p_1/2$，必有 $p_L > p_N$；消费者预见到被加价，第一期购买被抑制。}"杀熟"在理论上令人警惕的原因正在于此：它同时具备三级歧视的可行性与一级歧视的剩余榨取能力，而实施成本却低于两者。

行为价格歧视的争议集中在两点。其一，\textbf{公平性}：同等服务收取不同价格违背直觉的公平观念，且歧视对象（老客户、高活跃用户）往往是平台最忠实的用户，客观上构成对"忠诚"的惩罚。其二，\textbf{消费者剩余再分配}：个性化定价使剩余从高保留价格的消费者向平台转移，低保留价格的消费者则可能因更低报价而获益，总体福利方向与需求的异质程度有关；但"杀熟"针对的恰恰是弹性最低、最没有退路的那部分人，其分配后果尤其值得警惕。杀熟在消费者端的感知十分普遍：北京市消费者协会 2022 年的调查显示，约 82\% 的受访者认为大数据杀熟现象普遍存在，约 87\% 表示自己遭遇过差异化定价，高发场景依次为网购、在线旅游、外卖与网约车（北京市消费者协会，2022）。\mn{大数据杀熟的构成要件（论述题标准框架）：（1）数据：个体消费历史；（2）算法：推断支付意愿；（3）市场势力：缺乏替代选择的锁定用户。三者缺一则难以成立，这也是监管执法的查证难点：价格差异可能由成本差异或促销差异所致。}感知的普遍与认定的困难形成了鲜明反差，后者的成因在 5.2.4 节讨论。

\begin{misconception}
\textbf{误区一：价格不同就是杀熟。}机票、酒店、网约车的价格随时段与供需波动，属于\textbf{动态定价}（dynamic pricing）：定价依据是实时的供求状况，同一时刻查询的所有用户看到同一价格，与个人身份无关。判断的关键在于价格差异的依据，依据供需状况的是动态定价，依据个人特征的才是价格歧视。新用户优惠券、限时促销同样不构成杀熟，它们对所有符合条件者一视同仁。

\textbf{误区二：杀熟等同于一级价格歧视。}行为价格歧视所依据的是从历史行为推断出的支付意愿区间，精度远达不到个体保留价格的水平，榨取程度介于三级与一级之间。把它直接等同于完全价格歧视，会高估平台的信息能力，也会误判其福利后果。

\textbf{误区三：价格歧视本身违法。}价格歧视是普遍存在的定价策略，学生票、老年优惠、阶梯电价都属此列，法律并不一般性地禁止。受规制的是特定条件下的差别待遇：利用个人信息自动化决策实施的不合理差别待遇，以及具有市场支配地位的经营者对条件相同的交易相对人实行的差别待遇。合法与否取决于构成要件是否齐备，仅凭行为形态无法判断。
\end{misconception}

与个性化定价相对的另一侧现象是价格趋同。早期研究曾预期互联网会消除价格差异：搜寻成本趋近于零，消费者可以瞬间比价，同质商品的价格应当收敛到唯一水平。实证结果否定了这一预期，线上市场的价格离散（price dispersion）程度并不低于线下，同一商品在不同平台、不同时点的报价差异长期存在。\mn{线上价格离散三解释：（1）搜寻成本异质，比价者与不比价者并存，厂商在低价抢客与高价榨取之间取混合策略；（2）产品与服务差异，名义同质的商品实际附带不同的运费、时效与售后；（3）厂商定价策略，促销与个性化定价本身就在制造价差。}价格离散与本节主题直接相关：市场上本就存在大量价差时，消费者难以判断自己看到的报价究竟属于正常波动还是差别待遇，这正是杀熟举证困难的根源之一。

\subsection{个性化定价的福利分解}

5.2.1 节说明了平台如何获得定价所需的信息，接下来的问题是分配：个性化定价实施到极致时，福利在平台与消费者之间如何重新划分？这一分配可以精确地拆成两块，一块来自消费者剩余的转移，另一块来自无谓损失的转化。两块的性质截然不同，混在一起谈是福利评价中最常见的错误。以下用均匀分布的标准算例完成分解。

\begin{derivation}{个性化定价的福利分解}{}
设消费者的保留价格 $v$ 在 $[0,1]$ 上均匀分布，人口总量标准化为 1，平台的边际成本为零（数字服务的典型假设），每位消费者至多购买一单位。

\textbf{第一步，统一定价的基准。}平台对所有人收取同一价格 $p$，购买者为估价高于 $p$ 的消费者，人数为 $1-p$，利润为 $\pi_U = p(1-p)$。求导得 $1-2p=0$，故最优统一价格 $p^* = 1/2$，销量 $1/2$，利润
\[
\pi_U^* = \frac{1}{4}.
\]
消费者剩余等于购买者估价超出价格的部分之和，即 $(v-1/2)$ 在区间 $[1/2,1]$ 上的积分：
\[
CS_U = \int_{1/2}^{1}\Bigl(v - \frac{1}{2}\Bigr)dv = \Bigl[\frac{v^2}{2} - \frac{v}{2}\Bigr]_{1/2}^{1} = 0 - \Bigl(\frac{1}{8}-\frac{1}{4}\Bigr) = \frac{1}{8}.
\]
总剩余为 $TS_U = 1/4 + 1/8 = 3/8$。

\textbf{第二步，效率基准与无谓损失。}边际成本为零意味着任何一笔交易都创造正的社会价值，有效率的配置是全部人口成交，此时总剩余达到最大：
\[
TS^{\max} = \int_0^1 v\, dv = \frac{1}{2}.
\]
统一定价下的无谓损失即两者之差：
\[
DWL_U = \frac{1}{2} - \frac{3}{8} = \frac{1}{8}.
\]
它对应估价低于 $1/2$ 的那一半消费者，他们的支付意愿明明高于零边际成本，却被统一价格挡在了市场之外。

\textbf{第三步，完全个性化定价。}平台对每位消费者报出恰好等于其保留价格的价格 $p(v) = v$。此时所有人的净剩余均为零，按假设仍愿意成交，平台收入为
\[
\pi_P = \int_0^1 v\, dv = \frac{1}{2},
\]
消费者剩余 $CS_P = 0$，总剩余 $TS_P = 1/2$，无谓损失 $DWL_P = 0$。

\textbf{第四步，分解利润增量。}平台利润的增量为 $\Delta\pi = 1/2 - 1/4 = 1/4$。把它与前两步的结果对照，这一增量恰好由两部分构成：
\[
\underbrace{CS_U - CS_P = \frac{1}{8}}_{\text{原消费者剩余的转移}} \;+\; \underbrace{DWL_U - DWL_P = \frac{1}{8}}_{\text{原无谓损失的转化}} \;=\; \frac{1}{4}.
\]
两部分数量相同，性质却截然不同：\important{前一半是纯粹的再分配，消费者所失即平台所得；后一半是效率改善，原本不存在的交易被创造出来，总剩余因此增加}。
\end{derivation}

这一分解是评价个性化定价的标准框架，其价值在于把经常被混为一谈的两个问题分开。就总剩余而言，个性化定价是有益的：$TS$ 从 $3/8$ 升至 $1/2$，原本被高价挡在门外的低估价消费者获得了服务，交易范围扩展到全部人口。\mn{福利分解速记：$U[0,1]$、$MC=0$ 下，统一定价 $\pi=1/4$、$CS=1/8$、$DWL=1/8$；完全个性化 $\pi=1/2$、$CS=0$、$DWL=0$。利润增量 $1/4$ ＝ 剩余转移 $1/8$ ＋ 效率改善 $1/8$，各占一半。}就分配而言，个性化定价是消费者的净损失：$CS$ 从 $1/8$ 降至 $0$，全部剩余归于平台。福利评价的结论因此取决于所用的标准：以总剩余为准则，个性化定价改善福利；以消费者剩余为准则，个性化定价损害福利。

这一分歧也解释了监管态度的取向。中国的消费者权益保护与个人信息保护立法把差别待遇本身作为规制对象，其隐含的评价基准是消费者剩余，总剩余的增加不足以构成豁免理由。理由并不难理解：在垄断或强市场势力条件下，效率改善的收益全部被平台占有，消费者无从分享，两种基准的结论会系统性背离，此时选择消费者一侧的基准更符合立法的保护目的。

现实中的个性化定价达不到模型假设的完全榨取水平，分解框架依然适用。平台若只能识别一部分消费者，或估计存在误差，实际结果就落在统一定价与完全个性化之间：转移的规模按识别精度打折，效率改善的规模同样打折。评价具体案件时需要判断的是两块中哪一块占主导：以扩大覆盖为主的（如向价格敏感人群发放定向折扣以扩大用户基础），效率成分居多；以识别高支付意愿者并抬价为主的（杀熟的典型形态），转移成分居多。这一判断标准同样适用于下一小节的精准营销。

\subsection{精准营销的福利分析}

数据用于定价之外，还有一重更基础的用途是\textbf{精准匹配}：平台识别消费者的偏好类型，把合适的商品推给合适的人。精准营销同时改变两件事，匹配效率上升使消费者更容易找到心仪商品，细分能力增强又使厂商获得了差别定价的空间。两者的福利效果方向不同，一个数值算例可以完整展示这种张力。

\begin{derivation}{精准营销：匹配效应与定价效应的分离}{}
设市场上有黑、白两种产品，边际成本均为零。消费者共四类、每类一人：$H_B$ 对黑产品的估价为 100，$L_B$ 对黑产品的估价为 30，$H_W$ 与 $L_W$ 对白产品的估价对称地为 100 与 30。每位消费者只购买与自身偏好匹配的那种产品。

\textbf{情形一：随机匹配、统一定价（基准）。}厂商无法识别消费者类型，消费者也无法准确找到匹配自己偏好的产品，随机接触下只有一半的人遇到对的产品。先看定价：以黑产品为例，厂商面对 $H_B$（估价 100）与 $L_B$（估价 30）两类潜在买家，定价 100 时只有 $H_B$ 购买，定价 30 时两类都买，两种定价的收入分别为
\[
100 \times 1 = 100, \qquad 30 \times 2 = 60,
\]
故最优统一价格为 100。再计入匹配环节：只有一半消费者接触到对的产品，期望成交量减半，每种产品的期望收入为 50。两种产品合计
\[
\pi_1 = 50 \times 2 = 100, \qquad CS_1 = 0, \qquad TS_1 = 100.
\]
消费者剩余为零，因为唯一的成交者按其估价足额付款。

\textbf{情形二：精准匹配、统一定价。}平台用数据把产品推荐给偏好匹配的消费者，匹配率由 $1/2$ 提高到 $1$。定价规则不变，仍为统一价 100，但每种产品的期望成交量翻倍至 1 份：
\[
\pi_2 = 100 \times 2 = 200, \qquad CS_2 = 0, \qquad TS_2 = 200.
\]
与情形一相比总剩余增加 100，这一增量即\important{匹配效应}：它来自原本无法达成的交易被成功撮合，是新创造出来的价值。

\textbf{情形三：精准匹配、精准定价。}平台进一步识别消费者的估价类型，对 $H$ 类收 100、对 $L$ 类收 30，每种产品卖出 2 份：
\[
\pi_3 = (100 + 30) \times 2 = 260, \qquad CS_3 = 0, \qquad TS_3 = 260.
\]
与情形二相比总剩余增加 60，来自 $L_B$、$L_W$ 两位低估价消费者被纳入交易，这一增量即\important{定价效应}中的产出扩大成分。

\textbf{第四步，并列三种情形完成分解。}
\[
TS: 100 \to 200 \to 260, \qquad \pi: 100 \to 200 \to 260, \qquad CS: 0 \to 0 \to 0.
\]
总剩余的两次上升分别由匹配效应（$+100$）与定价效应（$+60$）贡献，而\important{消费者剩余在三种情形下恒为零，全部增量归于厂商}。
\end{derivation}

算例的关键结论在最后一行：总福利的扩大与消费者的处境是两回事。数据既提高了匹配效率（做大蛋糕），又强化了榨取能力（改变切法），在垄断且能够完全识别估价的极端假设下，两项收益全部落入厂商口袋，消费者的处境与数据出现之前完全相同。

这一结论对垄断假设高度敏感，竞争条件下结果会显著不同。多个平台争夺同一批用户时，匹配效率的提升会通过价格竞争部分传导给消费者；但精准投放同时使各厂商在自己的细分人群中更接近一个小垄断者，正面的价格竞争因此缓和。两种力量方向相反，净效果取决于细分市场之间的替代程度：细分越细、跨细分市场的替代越弱，竞争缓和的力量越强。

精准营销的福利结论因此可以概括为三句话：匹配效率上升扩大蛋糕，价格歧视改变切法，竞争强度决定切法向哪一方倾斜。\mn{精准营销福利算例速记：垄断下"匹配更多＋价格歧视"并存，总剩余随匹配率上升，消费者剩余看榨取程度；竞争下精准投放使各厂商成为细分市场的小垄断者，价格竞争缓和，总福利上升而消费者剩余方向不定。}对监管的含义十分明确：\important{评估数据驱动的营销行为，应当把匹配效应与定价效应分开计量}。只看价格歧视会低估数据的福利贡献，只看交易量上升又会掩盖剩余的转移。论述题按"总福利、消费者剩余、竞争条件"三层展开，即可覆盖完整的分析链条，第 6 章的竞争政策分析建立在同一框架之上。

\subsection{大数据杀熟的规制}

理论上的福利争议要落到执法，必须回答两个问题：现行法律依据是什么，认定的难点在哪里。中国对差别待遇的规制由数部法律法规从不同角度共同覆盖，分散的格局本身就反映了这一行为的复合属性：它既是价格行为，又是个人信息处理行为，还可能是滥用市场支配地位的行为。

最直接的依据是《个人信息保护法》（2021 年施行）第 24 条：个人信息处理者利用个人信息进行自动化决策，应当保证决策的透明度和结果公平、公正，不得对个人在交易价格等交易条件上实行不合理的差别待遇。这一条把规制落点放在"利用个人信息自动化决策"这一行为方式上，恰好对应行为价格歧视的技术实现路径。同法还赋予个人拒绝权：通过自动化决策方式作出对个人权益有重大影响的决定，个人有权要求说明并有权拒绝。《电子商务法》（2019 年施行）第 18 条则从推荐环节切入，要求经营者根据消费者的兴趣爱好、消费习惯等特征提供搜索结果的，应当同时提供不针对其个人特征的选项。该条针对的是个性化推荐，但推荐与定价的数据基础相同，"关闭个性化"的选项在实践中成为消费者对抗差别待遇的第一道防线。

规制链条的另外两环分别落在算法与价格披露上。《互联网信息服务算法推荐管理规定》（2022 年施行）第 21 条把禁止对象写得更为具体：算法推荐服务提供者向消费者销售商品或者提供服务的，不得根据消费者的偏好、交易习惯等特征，利用算法在交易价格等交易条件上实施不合理的差别待遇。\mn{杀熟规制法条速记：《个人信息保护法》第 24 条（自动化决策不得不合理差别待遇）、《电子商务法》第 18 条（须提供非个性化选项）、《算法推荐管理规定》第 21 条（不得依交易习惯差别定价）、《明码标价和禁止价格欺诈规定》（差别定价须标明条件）。}《明码标价和禁止价格欺诈规定》（2022 年施行）从价格披露角度加以约束：经营者对不同交易条件的交易对象实行不同价格的，应当标明相应的交易条件，隐瞒对消费者不利的价格条件可以构成价格欺诈。《反垄断法》的滥用市场支配地位条款也禁止没有正当理由对条件相同的交易相对人实行差别待遇，但其适用须先认定市场支配地位，门槛较高。

\textbf{《价格法》的适用争论}是本领域最值得注意的问题，也是辨析题的常见素材。《价格法》第 14 条禁止经营者"提供相同商品或者服务，对具有同等交易条件的其他经营者实行价格歧视"。条文的保护对象是\textbf{其他经营者}，规范的是经营者之间的横向歧视，例如供应商对不同经销商给出不同批发价，消费者并未被纳入该条的保护范围。这一取向源于立法时的经济背景：《价格法》制定于 1997 年，当时针对消费者的个体化定价在技术上尚不可能，立法者设想的价格歧视是流通环节的批发歧视。数字技术出现后，条文与现实之间形成了缺口，最典型、最普遍的价格歧视恰恰发生在经营者与消费者之间，而第 14 条无法直接适用于此。围绕这一缺口存在两种主张：一种主张修法扩展保护范围，理由是两类行为的经济实质相同，区别对待缺乏正当性；另一种主张维持现状，由《个人信息保护法》与《电子商务法》承担规制职能，理由是差别定价的合法性判断高度依赖个案，一般性禁止会误伤正常的促销与会员定价，而从个人信息处理角度切入更契合技术实现路径。

认定难点则集中在证明责任上。价格差异本身不构成违法，还须证明差异源于个人特征且不具有正当理由，而平台可以给出多种替代解释：新用户补贴、优惠券发放规则、渠道差异、时段与库存波动、会员等级权益。这些解释在形式上都属于"不同的交易条件"，定价算法的内部逻辑又属于商业秘密，消费者与监管者难以获取。举证困难与 5.2.1 节讨论的价格离散现象相互叠加，构成执法的实际障碍。近年的监管实践因此更多依赖平台自查、约谈与合规整改，较少依赖逐案的处罚认定。

\section{免费商业模式与注意力经济}

前两节讨论的定价工具都以收费为前提，区别只在收多少、向谁收。数字经济却把大量产品的标价定为零。免费属于价格结构的一种极端形态：某一侧的价格降到零甚至为负，成本由另一侧回收，这正是第 4 章倾斜式定价在商业层面的展开。

本节先辨析免费的三类逻辑（5.3.1 节），再深入 freemium 的版本切割设计（5.3.2 节），最后讨论免费模式最主要的收入来源，即注意力的出售与广告位的拍卖（5.3.3 节）。

\subsection{免费的三类逻辑}

免费是数字经济最普遍的定价形态，"免费"二字底下的经济学本质却各不相同。\mn{免费的三类逻辑速记：（1）交叉补贴：A 产品免费、B 产品收费；（2）三方市场：用户免费、广告主付费；（3）免费+增值（freemium）：基础版免费、高级版收费。}辨析商业模式时，只需回答一个问题：谁在为它付费。三种答案对应三类逻辑。

第一类是\textbf{交叉补贴}（cross-subsidization）：免费品与付费品在消费上互补，免费品充当付费品的营销投入，其亏损由付费品的利润覆盖。经典形态是低价打印机配高价墨盒、低价刀架配高价刀片，数字业务中的对应物是免费阅读器配收费电子书、免费基础工具配收费专业插件、开源框架免费而技术支持与托管服务收费。交叉补贴成立的条件是两种产品之间存在较强的互补性与一定的锁定：用户一旦购入打印机就被绑定在该品牌的耗材上，前期的免费投入才能在后续收回。锁定越弱，交叉补贴越容易失败，因为用户可以只拿走免费的那一半。

第二类是\textbf{三方市场}（three-sided market）：平台向用户免费提供内容与服务，同时向广告主出售用户的注意力，\important{免费用户本身构成被出售的"产品"}。这是双边市场倾斜式定价（第 4 章）最直接的应用：用户侧的价格为零甚至为负（免费叠加内容补贴与现金激励），广告侧承担全部成本并贡献利润。搜索引擎、社交平台、短视频、新闻资讯与多数手机游戏都属此类。倾斜的程度取决于两侧的弹性与交叉外部性强度：用户对价格极其敏感，收费会导致大规模流失；而用户规模越大，广告主的支付意愿越高。两个条件叠加，把用户侧价格压到零就成为最优选择。

第三类是 \textbf{freemium}：基础版本免费以获取用户基础，增值版本收费以回收成本，\important{免费版同时承担获客与价格锚定的双重功能}。与三方市场相比，freemium 的收入直接来自用户自身，免费用户与付费用户是同一批人中的两部分。云存储、在线协作工具、专业软件、音乐与视频会员多采用这一模式。它对产品形态有特定要求：产品必须能在功能维度上切分出层次分明的版本，切分点的位置直接决定转化率与收入，这正是下一小节要处理的问题。

三类逻辑的表层形态不同，经济学本质却一致：免费侧的亏损由另一侧的付费收入覆盖，免费是利润实现的一种结构安排。\important{免费不意味着没有收入，只意味着收入不来自使用行为本身}。付费方是同一用户在另一产品上（交叉补贴）、是广告主（三方市场）、还是同一用户在同一产品的高级版本上（freemium），决定了平台的全部经营逻辑：交叉补贴依赖锁定，三方市场依赖用户规模与注意力时长，freemium 依赖版本切割的精度。

\subsection{Freemium 的自选择设计}

三类免费逻辑中，freemium 的机制最为精细。平台无须给用户贴任何标签，只需设计好免费版与收费版的功能差异，让用户在版本选择中自行暴露其价值评价，这正是二级价格歧视的自选择思想在免费商业模式中的应用（5.1.2 节的自选择菜单是它的一般形式）。

\begin{definition}[freemium（免费增值模式）]
freemium 是 free 与 premium 的组合词，指基础版本免费、增值版本收费的商业模式。
\end{definition}

freemium 的设计核心是版本切割，也就是免费版保留多少功能。这个参数同时牵动两件事：它决定收费版的需求与定价空间，也决定免费版的获客能力。以下模型先刻画前一件事。

\begin{derivation}{Freemium 的版本切割与最优定价}{}
设消费者对产品完整功能的价值评价为 $v$，$v$ 在 $[0,1]$ 上均匀分布，人口总量标准化为 1，平台的边际成本为零。免费版提供的价值为 $v_F = \alpha v$，其中 $\alpha \in (0,1)$ 是版本切割参数，代表免费版保留的价值份额：$\alpha$ 越接近 1，免费版越接近完整功能；$\alpha$ 越接近 0，免费版越受限。收费版提供完整价值 $v$，定价为 $p$。

第一步，确定消费者的版本选择。选择收费版的净所得为 $v - p$，选择免费版的净所得为 $\alpha v$。消费者选择收费版当且仅当前者不低于后者：
\[
v - p \geq \alpha v.
\]
移项得 $v(1-\alpha) \geq p$，两边除以 $1-\alpha$（因 $\alpha<1$，该量为正，不等号方向不变），得到付费的临界价值
\[
v \geq \frac{p}{1-\alpha} \equiv \hat v.
\]
$\hat v$ 的表达式已经透露了核心机制：分母 $1-\alpha$ 是收费版相对于免费版的\textbf{增量价值系数}。免费版保留得越多（$\alpha$ 越大），付费所能换来的增量越小，临界值 $\hat v$ 越高，愿意付费的人越少。

第二步，求收费版的需求。均匀分布下，价值评价高于 $\hat v$ 的用户比例等于区间 $[\hat v, 1]$ 的长度，故收费版的需求为
\[
q(p) = 1 - \hat v = 1 - \frac{p}{1-\alpha}.
\]

第三步，求最优定价。平台利润等于价格乘以需求：
\[
\pi = p\Bigl(1 - \frac{p}{1-\alpha}\Bigr).
\]
这是关于 $p$ 的二次函数，开口向下，对 $p$ 求导得 $1 - 2p/(1-\alpha) = 0$，解出最优价格并代回利润函数：
\[
p^* = \frac{1-\alpha}{2}, \qquad \pi^* = \frac{1-\alpha}{4}.
\]
两式都是 $\alpha$ 的减函数，由此得到本模型的核心结论：\important{免费版保留的价值份额越小，收费版的最优价格越高、平台利润越大}。
\end{derivation}

推导给出的只是权衡的一半。若止步于此，平台应当把 $\alpha$ 压到尽可能小，让免费版几乎没有价值。现实中平台并不这样做，原因在于模型省略了免费版的另一重功能：获客。免费用户带来口碑传播、网络效应、数据积累与未来的转化机会，这些价值并不体现在当期的收费版利润里。把获客价值记为 $G(\alpha)$，它随 $\alpha$ 递增，因为免费版越好用，能吸引和留存的用户越多；平台的实际目标是最大化 $\pi^*(\alpha) + G(\alpha)$，最优切割点由两者的边际变化相等处决定。

版本切割点因此成为 freemium 设计的关键变量。切割过浅（$\alpha$ 接近 1），免费版已能满足绝大多数需求，无人付费；切割过深（$\alpha$ 接近 0），免费版形同虚设，用户基础无从积累，口碑与网络效应也无从发生。\mn{freemium 定价口诀：收费版最优价 $p^*=(1-\alpha)/2$、利润 $(1-\alpha)/4$；免费版保留份额 $\alpha$ 越小，收费版越贵、利润越高，但获客能力越弱。最优切割点落在"榨取"与"获客"的平衡处。}\important{最优切割点落在榨取能力与获客能力的平衡处}，这也是 freemium 与纯粹的二级价格歧视最大的区别：后者只需考虑剩余榨取，前者还要为用户规模留出空间。

现实中平台对切割维度的选择相当精细，而且往往同时使用多个维度。常见的切割维度包括容量（云存储的免费空间上限）、质量（视频清晰度、音乐码率）、广告（免费版插入广告，付费版去广告）、功能（协作工具的成员数上限、专业软件的高级功能）、时间（免费试用期）与授权范围（个人使用免费、商业使用收费）。多维切割的好处在于可以针对不同类型的用户分别设卡：轻度个人用户在容量维度上几乎不受限制，重度用户与团队用户则会先后撞上容量与成员数的上限。这实际上把 5.1.2 节的自选择菜单从一维推广到了多维，每个维度都是一道筛选。需要注意的是，freemium 的付费转化率在业内通常只有个位数百分比，收入因而高度依赖少数重度用户，产品设计必须同时服务两个目标：让大多数不付费的用户持续留存以维持网络效应与口碑，让少数高价值用户遇到足够明确的付费理由。

\subsection{注意力经济与广告竞价}

当用户免费而广告主付费时，平台出售的核心资源是\textbf{注意力}。注意力经济学的基本命题是：\important{注意力是稀缺资源}。信息的生产成本随数字技术持续下降，个体接收信息的带宽却受制于时间这一硬约束，信息过载与注意力稀缺因此成为一对孪生现象。西蒙（Herbert Simon）早在 1971 年即指出，信息的丰富制造注意力的贫困，稀缺的注意力由此获得经济价值，"注意力经济"（attention economy）一词亦源出于此。

注意力经济中的平台生产可以分解为两个阶段，安德森（Chris Anderson）在《免费》（2009）中把这一结构视为免费商业模式的核心。\important{第一阶段争夺注意力}：平台投入内容生产、创作者补贴与推荐算法，以免费服务换取用户的时间与关注，这一阶段只有成本、没有收入。\important{第二阶段出售注意力}：平台把汇集起来的注意力按人群、场景与时段打包，卖给广告主，收入在此实现。两个阶段的产品截然不同，前者的产品是内容与服务，后者的产品是用户的注意力；用户在前一阶段是需求方，在后一阶段却成为被交易的标的。这一结构解释了平台行为中若干看似矛盾的现象：平台不惜以巨额补贴争夺用户时长，因为时长是第二阶段的原料；平台持续优化推荐算法以延长使用时间，因为算法效率直接决定第一阶段的产出规模。中国短视频用户规模超过 10 亿（中国互联网络信息中心，2023），日均使用时长以小时计，注意力总量因此成为平台竞争最直接的标的，抖音、快手等平台的规模化运营正是这一逻辑的展开。

第二阶段的定价通过竞价拍卖完成，主流机制是广义第二价格拍卖（GSP）。

\begin{definition}[广义第二价格拍卖（generalized second-price auction，GSP）]
广义第二价格拍卖是搜索引擎广告位的拍卖机制：广告主对每次点击出价，广告位按出价（结合质量分）排序分配，获得第 $j$ 个位置的广告主支付第 $j+1$ 名的出价。
\end{definition}

早期搜索引擎曾使用按实际出价支付的广义第一价格拍卖，出价者为压低支付而反复试探降价，导致价格剧烈波动，谷歌与雅虎因而转向 GSP：出价与支付分离，竞价策略随之稳定。GSP 的名称容易让人以为它完整继承了第 2 章维克瑞拍卖的性质，实际情形要复杂一些，以下推导逐层澄清。

设广告位按点击率排序 $CTR_1 > CTR_2 > \cdots > CTR_K$，广告主 $i$ 对每次点击的估价为 $v_i$（私有信息），提交出价 $b_i$。图~\ref{fig:gsp} 以三个广告位、四个广告主为例，呈现配置与支付的结构。

\begin{figure}[H]
\centering
\begin{tikzpicture}[scale=0.85]
% 表格线：3 列（2.2 / 4.0 / 2.4）× 5 行
\draw[gray, thick] (0,0) rectangle (8.6,3.8);
\draw[gray] (2.2,0) -- (2.2,3.8);
\draw[gray] (6.2,0) -- (6.2,3.8);
\draw[gray] (0,3.0) -- (8.6,3.0);
\draw[gray] (0,2.25) -- (8.6,2.25);
\draw[gray] (0,1.5) -- (8.6,1.5);
\draw[gray] (0,0.75) -- (8.6,0.75);
% 表头
\node[font=\small\bfseries] at (1.1,3.43) {出价排序};
\node[font=\small\bfseries] at (4.2,3.43) {获得位置};
\node[font=\small\bfseries] at (7.4,3.43) {支付 $b_{j+1}$};
% 第 1 行
\node[font=\small] at (1.1,2.62) {第1名 $b_1$};
\node[font=\small] at (4.2,2.62) {位置1（点击率最高）};
\node[font=\small, color=winered] at (7.4,2.62) {$b_2$};
% 第 2 行
\node[font=\small] at (1.1,1.88) {第2名 $b_2$};
\node[font=\small] at (4.2,1.88) {位置2};
\node[font=\small, color=winered] at (7.4,1.88) {$b_3$};
% 第 3 行
\node[font=\small] at (1.1,1.12) {第3名 $b_3$};
\node[font=\small] at (4.2,1.12) {位置3};
\node[font=\small, color=winered] at (7.4,1.12) {$b_4$};
% 第 4 行
\node[font=\small, gray] at (1.1,0.38) {第4名 $b_4$};
\node[font=\small, gray] at (4.2,0.38) {落选};
\node[font=\small, gray] at (7.4,0.38) {—};
\end{tikzpicture}
\caption{GSP 的配置与支付}
\label{fig:gsp}
\end{figure}

这一规则的经济含义在于出价与支付的分离。在单一广告位的情形下，GSP 退化为第2章维克瑞拍卖：支付由他人的出价决定，出价只影响自己获得哪个位置，虚报出价只会扭曲自己获得的位置质量，因此\important{按真实估价出价} $b_i = v_i$ 是简单而可执行的策略；多位置情形下，获得位置 $j$ 的广告主为每次点击支付 $b_{j+1}$，竞价均衡中位置按估价排序配置（最高估价者获最佳位置），这是拍卖配置效率在广告场景的实现。GSP 的现实复杂性在于：竞价重复进行、估价动态变化、质量分（平台对广告质量的评分）介入排序，排名由"出价 × 质量分"决定，实际支付为维持该位置所需的最低出价，质量分高的广告主可以更低出价获得相同位置，平台由此在保持拍卖效率的同时兼顾用户体验。

GSP 机制的含义可以概括为两点：其一，出价与支付分离使诚实出价成为简单可执行的策略，位置配置在均衡中达到高效；\mn{GSP 记忆：搜索引擎广告位即广义第二价格拍卖；第 $j$ 位支付第 $j+1$ 名的出价；单一广告位下占优策略是按真实点击估价出价；广告平台"用户免费"的本质是用户以注意力支付。}其二，注意力变现的完整链条由此闭合：平台免费获取用户注意力，通过拍卖将其转化为广告收入，\important{用户以注意力支付了"免费"服务}。

\section{平台佣金与抽成}

\subsection{佣金的经济学}

平台佣金（take rate）是平台向交易收取的成交额比例，是交易平台的核心收入模式。佣金本质上是平台对价格结构的再设计：消费者支付的价格中多大比例流向商家、多大比例归属平台，是双边市场价格结构规则（第4章）在收入端的直接体现。\mn{佣金归宿速记：弹性更小的一侧承担更多；商家对平台依赖强（弹性小）则"抽商家"，消费者对价格敏感（弹性大）则平台自行吸收。}佣金率的经济效应可以通过双边市场框架分析。

\begin{derivation}{佣金归宿与最优佣金率}{}
考虑一个交易平台，设消费者支付价格 $p_B$，商家获得净价格 $p_S$，佣金率 $t$（对成交额征收）满足 $p_S = p_B(1-t)$；消费者需求的价格弹性为 $\eta_B = -\frac{p_B}{D}\frac{dD}{dp_B}$，商家供给的价格弹性为 $\eta_S = \frac{p_S}{S}\frac{dS}{dp_S}$（均取绝对值）。佣金等价于对交易征收的从价税，其归宿由税负转嫁分析给出：对市场均衡条件 $D(p_B) = S(p_B(1-t))$ 两边关于 $t$ 全微分，代入弹性定义整理，得
\[
\text{商家承担份额} = \frac{\eta_B}{\eta_S + \eta_B}, \qquad
\text{消费者承担份额} = \frac{\eta_S}{\eta_S + \eta_B},
\]
即\important{弹性更小的一侧承担更多佣金}（税负归宿的一般结论）：商家侧弹性小（对平台依赖强）时，佣金的大部分由商家承担，表现为商家到手的净价格下降；消费者侧弹性大时，消费者承担的份额小，平台不敢将佣金转嫁给消费者。这解释了外卖、出行平台"抽商家不抽消费者"的佣金设计：商家侧的依赖性与弹性结构决定了佣金归宿。

\textbf{最优佣金率}：平台提供交易基础设施的边际成本近似为零，利润为 $t \cdot V(t)$（$V$ 为交易量），一阶条件给出
\[
\frac{d\pi}{dt} = V + t\,V' = 0 \quad\Rightarrow\quad \varepsilon_V(t^*) \equiv -\frac{t^*\,V'(t^*)}{V(t^*)} = 1,
\]
其中 $\varepsilon_V$ 为交易量对佣金的弹性（绝对值）。这是勒纳公式在零边际成本下的特例：\important{最优佣金率应使交易量弹性恰为 1}。佣金率提高 1\% 使交易量下降超过 1\% 时，佣金率已经过高。交易量对佣金越敏感（商家替代渠道越多、消费者对价格越敏感），弹性升至 1 所需的佣金率越低，最优佣金率越低。
\end{derivation}

\subsection{App Store 30\% 抽成之争}

苹果与谷歌对应用商店内的数字内容交易收取 15\%–30\% 的佣金，这一抽成比例的争议是平台佣金问题的标志性案例。\mn{案例：Epic 诉苹果案（2020–2021）：Epic 在《堡垒之夜》中绕过苹果支付渠道，触发下架与诉讼。法院最终裁定苹果的抽成不构成垄断，但禁止其"反引导"条款（阻止开发者告知用户其他支付渠道）。}从双边市场理论看，应用商店是典型的软件平台：开发者侧单归属倾向强（iOS 用户只在 App Store 装机），用户侧被锁定于硬件生态。30\% 抽成的可持续性恰恰源于这一锁定结构，其数量关系可以用竞争瓶颈框架核算。

\begin{derivation}{抽成上限的核算}{}
设开发者的应用在 App Store 年收入为 $R$，抽成率 $t$，开发者的替代选择（放弃 iOS、仅上架安卓）带来的收入损失比例为 $\ell$（即 iOS 收入占比），则开发者"留用与退出无差异"的条件为
\[
R(1 - t) = R(1 - \ell) \;\Rightarrow\; t_{\max} = \ell.
\]
抽成率的上限即\important{被锁定收入的比例}：开发者的 iOS 收入占比越高，平台可提取的佣金上限越高。对重度依赖 iOS 的开发者，$\ell$ 接近 $1$，理论上可承受的抽成远超 30\%——30\% 的争议性不在于它逼近了开发者退出的临界点，而在于其\textbf{形成过程}：它不是通过竞争性谈判形成，而是平台单方设定、开发者"接受或退出"。这一区别的经济学实质是竞争瓶颈（第4章）的定价性质：单归属侧的价格不是竞争均衡价格，而是垄断性收费，其水平由锁定程度（$\ell$）而非成本或竞争约束决定。App Store 的辩护逻辑（基础设施投资回收、安全审核的公共品供给）在理论上不改变这一结论，只提供效率抗辩，争议的裁决取决于平台公共品投资与垄断收费之间的权衡，这正是 Epic 案判决"不构成垄断、但禁止反引导"的折中逻辑。
\end{derivation}

\subsection{接入费与交易费的组合}

平台的收费工具不止佣金一种：\textbf{接入费}（access fee，会员费、入驻费）对加入收费，\textbf{交易费}（transaction fee，佣金）对交易量收费。最优收费结构的选择取决于两种费用的比较静态。

\begin{derivation}{接入费与交易费}{}
接入费 $A$ 与交易费 $t$ 的组合利润为
\[
\pi = A \cdot N(A, t) + t \cdot V(A, t),
\]
其中 $N$ 为用户规模、$V$ 为交易量。接入费在用户参与决策前征收，抑制参与但对交易无扭曲（准入门槛）；交易费按交易量征收，扭曲交易激励但随交易规模自动扩缩收入。最优结构的一般结论：\important{交易不确定性强、用户间交易量差异大时，交易费优于接入费}（按使用量收费而非一刀切）；\important{平台需要快速建立用户基础时，接入费应低甚至为零}（以交易费为主要回收渠道）。现实中平台普遍以"免费接入 + 交易抽成"起步，成熟后叠加会员制（接入费），正是对两种工具性质的实践运用。
\end{derivation}

\begin{chapterbox}
\textbf{本章考点清单}

\begin{enumerate}
\item 三级价格歧视：$MR_1 = MR_2 = MC$ 推导，弹性小的群体价格高 \important{【专硕·核心】【学硕·热点】}
\item 两部收费制：$p^* = c$、入场费榨取剩余 \important{【专硕·核心】}
\item 一级价格歧视：产出效率与完全竞争相同、分配后果 \important{【专硕·核心】}
\item 行为价格歧视两期模型：老客户被加价的推导 \important{【专硕·核心】}
\item 大数据杀熟的构成要件与规制依据 \important{【专硕·核心】}
\item 免费的三类逻辑与 freemium 版本切割（$p^* = (1-\alpha)/2$）\important{【专硕·核心】}
\item 注意力经济与 GSP 广告竞价 \important{【专硕·扩展】}
\item 精准营销的福利算例：匹配效应扩大总福利、定价效应重塑分配 \important{【专硕·扩展】}
\item 佣金税负归宿与最优佣金率 $\varepsilon_V$ 规则 \important{【专硕·核心】}
\item App Store 抽成争议的竞争瓶颈分析 \important{【专硕·扩展】}
\item 接入费与交易费的最优组合 \important{【专硕·扩展】}
\end{enumerate}
\end{chapterbox}

本章完成了平台"如何收费"的完整图景：从价格歧视的经典工具，到大数据时代的个性化定价，再到免费与佣金的具体商业形态。值得注意的是，本章的每一种定价手段——歧视、免费、抽成——在越过某一界限后都同时是竞争法审查的对象；定价与竞争的规范边界，留待下一章处理。
